\documentclass{article}
\usepackage{graphicx} % Required for inserting images
\usepackage{amsfonts}

\title{Resoluçaõ dos Exercícios 4 e 5, Capítulo 6

Stein Shakarchi, Fourier Analysis}
\author{Afonso Francisco }
\date{February 2026}

\begin{document}

\maketitle

\textbf{4a)}
Usando a sugestão:
$$1=
\int_{\mathbb{R}^d} e^{-\pi |x|^2}dx=
\int_{S^{d-1}}\int_0^\infty e^{-\pi r^2}r^{d-1}drd\sigma(\gamma)=
\int_{S^{d-1}}d\sigma(\gamma) \int_0^\infty e^{-\pi r^2}r^{d-1}dr$$
$$=A_d \int_0^\infty e^{-\pi r^2}r^{d-1}dr$$

Usando a mudança de variáveis: $t=\pi r^2$ temos: $r=(\frac{t}{\pi})^{1/2}$ e $dr=\frac{dt}{2(\pi t)^{1/2}}$. Portanto,

$$1=
A_d\int_0^\infty e^{-t} \left( \frac{t}{\pi} \right)^{\frac{d-1}{2}} \frac{1}{2(\pi t)^{1/2}} dt=
\frac{A_d}{2\pi^{d/2}}\int_0^\infty e^{-t}t^{\frac{d}{2}-1}dt=
\frac{A_d}{2\pi^{d/2}}\Gamma\left(\frac{d}{2}\right)$$

Portanto,

$$A_d=\frac{2\pi^{\frac{d}{2}}}{\Gamma\left( d/2 \right)}$$\\

\textbf{b)}
Seja $B$ a bola unitária em $\mathbb{R}^d$. Temos,
$$V_d=
\int_BdV=
\int_{S^{d-1}}\int_0^1 r^{d-1}drd\sigma(\gamma)=A_d \left[ \frac{r^d}{d} \right]_0^1=\frac{A_d}{d}$$

Portanto, $$dV_d=A_d$$

Logo $$V_d=\frac{2\pi^{\frac{d}{2}}}{d\Gamma\left( d/2 \right)}=
\frac{\pi^{\frac{d}{2}}}{\frac{d}{2}\Gamma\left( d/2 \right)}=
\frac{\pi^{\frac{d}{2}}}{\Gamma\left( d/2+1 \right)},$$

onde na última igualdade usámos a propriedade da função gamma: $\Gamma(z+1)=z\Gamma(z)$.\\


\textbf{5.} Sendo $A$ simétrica, pelo \textit{Teorema espectral} vem que podemos escrever 
$$A=RDR^{-1},$$
onde $R$ é rotação ($R^T=R^{-1}$ e $|det \;R|=1$) e $D$ diagonal com entradas $\lambda_1,...,\lambda_d$, onde $\{\lambda_i\}$ são os valores próprios de $A$
\footnote{O teorema espectral diz que $A=QDQ^T$, onde $Q$ é ortogonal com colunas correspondendo aos vetores próprios de A.
Se $\det(Q)=-1$, podemos multiplicar um dos vetores por $-1$ fazendo com que $\det(Q)=1$. No entanto, o novo vetor continua a ser um vetor próprio colinear com o original, o que não altera a ortogonalidade nem alterna $D$. Assim, podemos sempre assumir que $Q=R$ é uma rotação.}.

Temos 
$$(x,A(x))=
x^TAx=x^TRDR^{-1}x=
x^TRDR^Tx$$

Usando a mudança de variávies $y=R^Tx$, temos $x=Ry$ e $dx=|det \;R|dy=dy$, pois $det\;R=1$, logo
$$(x,A(x))=
x^TRDR^Tx=
y^TDy=
\sum_{i=1}^d\lambda_iy_i^2,$$

Portanto,

$$\int_{\mathbb{R}^d}e^{-\pi(x,A(x))}dx=
\int_{\mathbb{R}^d}e^{-\pi\sum_{i=1}^d\lambda_iy_i^2}dy=
\int_{\mathbb{R}^d} \prod_{i=1}^d e^{-\pi\lambda_iy_i^2}dy=
\prod_{i=1}^d \int_{\mathbb{R}} e^{-\pi\lambda_iy_i^2}dy_i=$$
$$\prod_{i=1}^d \lambda_i^{-1/2}=
\left( \prod_{i=1}^d \lambda_i \right)^{-1/2}=
(det\; A)^{-1/2},$$

Onde na quarta igualdade usámos a mudança de variáveis $x_i=\lambda_i^{1/2}y_i$, transformando-o no integral $\lambda_i^{-1/2}\int_{\mathbb{R}}e^{-\pi x_i^2}dx_i = \lambda_i^{-1/2} \\


A sexta igualdade é justificada pelo seguinte argumento:
Seja $A$ uma matriz $d \times d$, fatorizando o polinómio característico vem 
$$p(\lambda)=\det(\lambda I-A)=(\lambda-\lambda_1)...(\lambda-\lambda_d)$$
Fazendo $\lambda =0$, 
$$p(0)=\det(-A)=(-\lambda_1)...(-\lambda_d)=(-1)^d \prod_{i=1}^d\lambda_i$$
Finalmente, já que $\det(-A)=(-1)^d \det(A)$, temos 
$$\det(A)= \prod_{i=1}^d\lambda_i$$






\end{document}
